% Supplementary A — Conference paper §I (Introduction)
% Extended motivation, historical context, and exhaustive prior-work
% citations that did not fit in the 4-page conference manuscript.

\documentclass[11pt,a4paper]{article}
% Shared preamble for all supplementary chapter documents.
% Used by each supplementary/conference/suppl_*.tex and
% supplementary/journal/suppl_*.tex via % Shared preamble for all supplementary chapter documents.
% Used by each supplementary/conference/suppl_*.tex and
% supplementary/journal/suppl_*.tex via % Shared preamble for all supplementary chapter documents.
% Used by each supplementary/conference/suppl_*.tex and
% supplementary/journal/suppl_*.tex via \input{../preamble.tex}.

\usepackage[utf8]{inputenc}
\usepackage[T1]{fontenc}
\usepackage[a4paper,margin=2.2cm]{geometry}
\usepackage{amsmath,amssymb,amsfonts,amsthm}
\usepackage{mathtools}
\usepackage{graphicx}
\usepackage{booktabs}
\usepackage{array}
\usepackage{multirow}
\usepackage{longtable}
\usepackage{algorithm}
\usepackage{algpseudocode}
\usepackage{xcolor}
\usepackage{url}
\usepackage{cite}
\usepackage[hidelinks]{hyperref}

\graphicspath{{../../figures/}}

\renewcommand{\thesection}{\Alph{section}.\arabic{section}}
\renewcommand{\thesubsection}{\thesection.\arabic{subsection}}

\theoremstyle{plain}
\newtheorem{theorem}{Theorem}[section]
\newtheorem{lemma}[theorem]{Lemma}
\newtheorem{proposition}[theorem]{Proposition}
\theoremstyle{definition}
\newtheorem{definition}[theorem]{Definition}
\theoremstyle{remark}
\newtheorem{remark}[theorem]{Remark}

\setcounter{secnumdepth}{3}
\setcounter{tocdepth}{2}

\newcommand{\R}{\mathbb{R}}
\newcommand{\E}{\mathbb{E}}
\newcommand{\KL}[2]{\mathrm{KL}\!\left(#1 \,\middle\|\, #2\right)}
\newcommand{\ari}{\mathrm{ARI}}
\newcommand{\nmi}{\mathrm{NMI}}
\newcommand{\softmax}{\mathrm{softmax}}
.

\usepackage[utf8]{inputenc}
\usepackage[T1]{fontenc}
\usepackage[a4paper,margin=2.2cm]{geometry}
\usepackage{amsmath,amssymb,amsfonts,amsthm}
\usepackage{mathtools}
\usepackage{graphicx}
\usepackage{booktabs}
\usepackage{array}
\usepackage{multirow}
\usepackage{longtable}
\usepackage{algorithm}
\usepackage{algpseudocode}
\usepackage{xcolor}
\usepackage{url}
\usepackage{cite}
\usepackage[hidelinks]{hyperref}

\graphicspath{{../../figures/}}

\renewcommand{\thesection}{\Alph{section}.\arabic{section}}
\renewcommand{\thesubsection}{\thesection.\arabic{subsection}}

\theoremstyle{plain}
\newtheorem{theorem}{Theorem}[section]
\newtheorem{lemma}[theorem]{Lemma}
\newtheorem{proposition}[theorem]{Proposition}
\theoremstyle{definition}
\newtheorem{definition}[theorem]{Definition}
\theoremstyle{remark}
\newtheorem{remark}[theorem]{Remark}

\setcounter{secnumdepth}{3}
\setcounter{tocdepth}{2}

\newcommand{\R}{\mathbb{R}}
\newcommand{\E}{\mathbb{E}}
\newcommand{\KL}[2]{\mathrm{KL}\!\left(#1 \,\middle\|\, #2\right)}
\newcommand{\ari}{\mathrm{ARI}}
\newcommand{\nmi}{\mathrm{NMI}}
\newcommand{\softmax}{\mathrm{softmax}}
.

\usepackage[utf8]{inputenc}
\usepackage[T1]{fontenc}
\usepackage[a4paper,margin=2.2cm]{geometry}
\usepackage{amsmath,amssymb,amsfonts,amsthm}
\usepackage{mathtools}
\usepackage{graphicx}
\usepackage{booktabs}
\usepackage{array}
\usepackage{multirow}
\usepackage{longtable}
\usepackage{algorithm}
\usepackage{algpseudocode}
\usepackage{xcolor}
\usepackage{url}
\usepackage{cite}
\usepackage[hidelinks]{hyperref}

\graphicspath{{../../figures/}}

\renewcommand{\thesection}{\Alph{section}.\arabic{section}}
\renewcommand{\thesubsection}{\thesection.\arabic{subsection}}

\theoremstyle{plain}
\newtheorem{theorem}{Theorem}[section]
\newtheorem{lemma}[theorem]{Lemma}
\newtheorem{proposition}[theorem]{Proposition}
\theoremstyle{definition}
\newtheorem{definition}[theorem]{Definition}
\theoremstyle{remark}
\newtheorem{remark}[theorem]{Remark}

\setcounter{secnumdepth}{3}
\setcounter{tocdepth}{2}

\newcommand{\R}{\mathbb{R}}
\newcommand{\E}{\mathbb{E}}
\newcommand{\KL}[2]{\mathrm{KL}\!\left(#1 \,\middle\|\, #2\right)}
\newcommand{\ari}{\mathrm{ARI}}
\newcommand{\nmi}{\mathrm{NMI}}
\newcommand{\softmax}{\mathrm{softmax}}


\title{Supplementary A --- Extended motivation for the band-mask
robustness diagnostic\\
{\large Companion to: \emph{A Band-Mask Robustness Diagnostic for LDA on HSI}
(WHISPERS 2026)}}
\author{Felipe Santib\'a\~nez-Leal}
\date{2026}

\begin{document}
\maketitle

\section{Scope}
This supplement expands Section~I of the conference manuscript. The
main paper limits itself to a single contribution --- a band-mask
robustness diagnostic for LDA on HSI --- presented in four printed
pages. Here we provide (a)~the historical context that motivated the
diagnostic, (b)~a complete enumeration of the failure modes the
diagnostic intends to surface, (c)~a contrast with adjacent
diagnostics in the literature, and (d)~the threats to validity we
acknowledge upfront.

\section{Historical context}

Latent Dirichlet Allocation~\cite{blei2003lda} introduced the
generative-topic-model framing of a discrete-token corpus over a
decade before its first reported application to hyperspectral
imagery. The application to HSI~\cite{zou2017pmlda,liu2019hsilda} was
preceded by two adjacent paradigms whose evaluation methodology
shaped the field's expectations:
linear spectral unmixing~\cite{nascimento2005vca,bioucasdias2012unmixing}
and deep spectral--spatial
classification~\cite{chen2014deep,paoletti2019deep}. Both paradigms
matured under an \emph{accuracy-first} evaluation discipline: the
canonical figure of merit is overall accuracy (OA) of a downstream
classifier on a held-out fold of one of the standard labelled
benchmarks (Indian Pines, Salinas, Pavia~U; later KSC and Botswana),
and methodological progress is measured as a delta against the best
prior OA on the same scene.

This discipline served unmixing and deep classification well: the
inductive bias of each paradigm is sufficiently strong that
post-hoc interpretability claims (endmember spectra; feature
visualisations) carry low risk of being a freak of a particular
training run. Topic models do not enjoy that protection. The basis
$\phi \in \R_+^{K \times B}$ inferred by LDA is identifiable only up
to a label permutation, and the inference procedure (online
variational Bayes~\cite{hoffman2010online}, collapsed Gibbs
sampling~\cite{griffiths2004finding}, or amortised variational
autoencoders in the neural variants~\cite{srivastava2017autoencoding,
dieng2020topic}) has multiple posterior modes for any finite-data
problem. The interpretability of a single trained model is therefore
contingent on a chain of stability properties:

\begin{enumerate}
\item \textbf{Seed stability}: refitting under a different random
initialisation should yield the same basis up to permutation.
\item \textbf{Capacity stability}: increasing $K$ by one should split
one canonical topic into two, not reshuffle every topic.
\item \textbf{Sample stability}: subsampling the corpus should not
relabel topics.
\item \textbf{Hyperparameter stability}: nudging $\alpha$ or $\eta$
should not alter topic identity at the same $K$.
\item \textbf{Spectral-window stability}: dropping the
atmospheric-water bands, or restricting to VNIR or SWIR only, should
not relabel topics (modulo the bands that the mask actually removes).
\end{enumerate}

Failure of any one of these properties weakens any claim of the form
``topic~$k$ is a mineral signature'' or ``topic~$k$ corresponds to
healthy canopy''. The first four properties are reported (sparingly)
in the existing literature; the fifth --- spectral-window stability
--- is, to our knowledge, not reported anywhere as a per-scene
diagnostic. The band-mask diagnostic of the main paper fills exactly
this gap.

\section{Failure modes the diagnostic surfaces}

The band-mask diagnostic is a hypothesis test under the null:
\textsc{H$_0$}: \emph{The canonical topic basis is a property of the
scene's underlying material distribution and is therefore robust to
the choice of spectral window.}

The diagnostic returns a per-(scene, mask) paired ARI; values close
to 1 fail to reject \textsc{H$_0$}, values close to 0 reject. There
are three substantively different ways to reject:

\begin{itemize}
\item \textbf{Mode (i): canonical-fit redundancy.} The canonical
basis exploits fine-grained correlations that are present in the
full-band representation and disappear under restriction. The
canonical topics are correct in the canonical reference frame but
not transportable across spectral windows.
\item \textbf{Mode (ii): weak identification.} The latent basis is
under-determined by the available data (small per-class sample, large
class count, narrow bands). Different posterior modes are equally
plausible under different band subsets.
\item \textbf{Mode (iii): mask-induced collinearity.} The mask
removes the bands that were carrying the discriminative signal, so
the masked LDA fits a different basis on the residual spectrum that
is correct for the residual but incomparable to the canonical.
\end{itemize}

The diagnostic does not separate (i)--(iii) by design --- doing so
would require additional information (e.g.\ band-wise Fisher ratios,
sample-size posterior bounds) and would inflate the diagnostic's
artefact footprint without changing the per-scene yes/no judgement
that motivates it. The conference paper's discussion explicitly flags
that the diagnostic is a \emph{flag-raising} tool, not a
\emph{root-cause-isolating} tool.

\section{Contrast with adjacent diagnostics}

We briefly contrast the band-mask diagnostic with three related
diagnostics in the wider topic-model and clustering literature.

\paragraph{Seed-stability swarms.} Plotting per-seed ARI to ground
truth across $N$ refits~\cite{koltchinskii2000empirical,
vinh2010information} answers ``does the same model on the same data
yield the same topics'', a different and weaker question than ``does
the same model on related data yield the same topics''. Seed
stability is, in our framework, axis F-3; the band-mask diagnostic is
axis F-5.

\paragraph{Cross-fold ARI.} K-fold CV on $\theta$ followed by ARI
comparison is sensitive to the held-out fold's class proportion; it
does not surface whether the topic basis itself is stable.

\paragraph{Coherence measures.} $c_v$ and $c_{\text{NPMI}}$
\cite{roder2015exploring} score the top-$N$ words per topic against
a reference corpus; they capture intra-topic compactness but say
nothing about cross-fit identity preservation.

The band-mask diagnostic occupies a different methodological niche
from each of the above: it asks whether the inferred topic
\emph{identities} survive a specific class of data restriction, with
the answer interpretable directly as a chance-corrected proportion.

\section{Threats to validity acknowledged upfront}

\paragraph{Top-50 Fisher ranking is computed on the unmasked
spectrum.} The Fisher ratio per band is computed against the
ground-truth label on the full-band spectrum; the mask then selects
the top-50 bands by that ranking. This is a deliberate choice ---
recomputing the Fisher ranking on each masked subset would conflate
the within-mask discriminability with the between-mask comparison
the diagnostic is trying to surface. The trade-off is that the
top-50 Fisher mask is a \emph{post-hoc} selector and its paired-ARI
value should be interpreted accordingly.

\paragraph{Canonical $K$ is held fixed.} The mask-refit uses the
canonical $K$ rather than a per-mask $K^*$. We argue this is the
right comparison because changing $K$ between fits introduces a
second degree of freedom that mixes with the band-restriction we are
isolating. The cost is that on a scene where the canonical $K$ is too
large for the masked spectrum, the masked fit will under-utilise its
capacity.

\paragraph{Sentinel masking of unlabelled pixels.} Both the canonical
and masked dominant-topic maps use sentinel 255 for unlabelled pixels;
the paired ARI is computed only on the intersection of labelled
pixels. The two fits therefore see a different document set from the
band axis but the same document set from the spatial axis.

\bibliographystyle{IEEEtran}
\bibliography{../../bibliography/refs}

\end{document}
